%%%%%%%%%%%%%%%%%%%%%%%%%%%%%%%%%%%%%%%%%%%%%%%%%%%%%%%%%%%%%%%%%%%%%%%%%%%%%%%%
%% 
\cleardoubleoddpage%  Make sure to start each chapter on a new odd page
\chapter{Background}
\label{sec:background}


\section{Fundamentals of Reverse Engineering}
\label{sec:background:fundamentals}

Reverse engineering aims to understand the internal behavior and structure of a program in the absence of its source code. 
This can be achieved either by closely observing a program during execution or, as demonstrated in this thesis, by statically analyzing the compiled binary itself.

In the case of compiled software, analysis begins with machine code, which represents the lowest level of abstraction. 
This machine code can be translated into assembly instructions, making the program logic more accessible to human analysts, as directly interpreting 
raw hexadecimal representations is impractical. Some analysis tools go a step further and attempt to reconstruct assembly instructions into higher-level source code, 
typically resembling the C programming language. While this reconstructed code is rarely identical to the original source, it often provides a functional approximation 
of the program's logic. The limitations of this process stem from compiler optimizations, which can significantly alter control flow and data structures, as well as 
from the fact that not all programs are originally written in C.

This thesis focuses on the static analysis of malware. In this context, reverse engineering is typically an expected and deliberately hindered activity. 
Malware authors frequently employ techniques designed to obstruct analysis and conceal malicious behavior. These techniques include, but are not limited to, 
control-flow obfuscation, encrypted execution stages, delayed payload execution, communication with remote command-and-control servers, and the use of executable packers. 
These methods are addressed in more detail later in this work.

Despite these challenges, understanding malware behavior is of critical importance. Identifying persistence mechanisms enables effective removal of infections, 
while analyzing command-and-control communication is essential when investigating botnets. Furthermore, reconstructing payload logic supports incident response efforts 
and can provide valuable insights for subsequent forensic investigations.

Not all malware exists in the form of native binaries. Some malicious software is distributed as scripts, while others target virtual machines or intermediary languages. 
Although this thesis primarily focuses on the analysis of compiled binaries, tools for analyzing Java and C\# programs are also included to account 
for these common malware formats.

\section{Static vs. Dynamic Analysis}
\label{sec:background:staticVsDynamic}

Dynamic analysis refers to the examination of malware by executing either parts of the program or the entire binary within a controlled environment. 
During execution, the analyst observes the runtime behavior of the sample. Typical aspects of interest include file system interactions, created or modified processes, 
registry changes, network communication, and other system level side effects. Specialized tools support this process by monitoring system calls, tracking file handles, 
recording process activity, and capturing network traffic.

In contrast, static analysis is performed without executing the program. Instead, the binary itself is examined directly. This may involve disassembly, decompilation, string analysis, 
control flow reconstruction, and inspection of embedded resources. The goal is to understand the program's structure and logic solely through examination of its code and metadata.

Both approaches have distinct advantages. Dynamic analysis can produce rapid and practical insights into malware behavior, particularly when dealing with heavily obfuscated or encrypted binaries, 
since runtime execution naturally reveals decrypted payloads and unpacked code. Furthermore, observing network communication during execution often provides valuable intelligence regarding 
command-and-control infrastructure and data exfiltration mechanisms.

However, sophisticated malware frequently incorporates anti-analysis techniques designed to detect virtualized or sandboxed environments. Such context aware behavior may alter or suppress malicious 
actions when analysis is suspected, thereby limiting the reliability of dynamic observations. Static analysis is not susceptible to this specific form of evasion, as no execution takes place. 
Consequently, it provides a more controlled and deterministic perspective on the program's structure, although it may require significantly more effort when strong obfuscation techniques are employed.


%% 
%%%%%%%%%%%%%%%%%%%%%%%%%%%%%%%%%%%%%%%%%%%%%%%%%%%%%%%%%%%%%%%%%%%%%%%%%%%%%%%%
